\section{Discussion}
\subsection{Interpretation of Results}
The primary objective of this study was to evaluate the usability, perceived usefulness, and cognitive workload of a student-facing JupyterLab dashboard designed to visualize fine-grained notebook interactions. The results provide a differentiated perspective on how students perceive learner-facing analytics tools in programming contexts and allow the findings to be situated within the broader literature on notebook analytics and learning dashboards.
\subsubsection{RQ 1: Usability}
The mean SUS score of $M=65.57$ indicates that the system is marginally acceptable according to established benchmarks. While this score falls just short of the above average threshold of 68, it is important to contextualize this within the internal consistency of the scale $\alpha=.794$. The gap between the observed score and the benchmark suggests that while the core functionality is present, the integration of complex interaction data, such as cell execution histories and error patterns, into a single dashboard may require further UI refinement to reduce perceived complexity. Qualitative feedback from participants explicitly illustrated these interface constraints. For instance, one student suggested that ``in top error cells it might be better if it could show error lines in the code and in the dashboard.'' This highlights a clear user desire for tighter visual integration between the analytics and the active coding area. Furthermore, one potential factor contributing to the SUS score falling slightly below the average threshold relates to the dashboard's layout and information density. Due to limited screen real estate, participants were required to scroll vertically to access all visualizations. This necessity for scrolling and searching for information can increase the cognitive burden on users, a challenge that learning analytics tools must actively avoid to be effective~\cite{corrin2014exploring}. Transitioning to a single-page layout that presents all summary statistics and timelines simultaneously would align better with the established design principle of perceiving learning at a glance~\cite{schwendimann2016perceiving}. Such a redesign could enhance clarity, reduce navigational effort, and potentially lead to higher technology acceptance and usability ratings.

\subsubsection{RQ 2: Usefulness and Ease of Use}
A compelling contrast emerged between the TAM subscales. Participants rated the Perceived Ease of Use (PEOU) significantly higher $M = 5.78$ than the Perceived Usefulness (PU) $M = 4.93$. The excellent reliability of the PEOU scale $\alpha = .941$ underscores that the dashboard itself is perceived as intuitive and easy to navigate. These findings strongly align with existing literature on learning analytics. The high PEOU score demonstrates that the dashboard successfully meets the requirement outlined by Corrin and De Barba~\cite{corrin2014exploring}, who emphasize that reflection tools must be seamlessly integrated without increasing cognitive burden. However, the more moderate PU score reflects a well-documented challenge in the field as Corrin and De Barba~\cite{corrin2014exploring} also note, students often require time and guidance to effectively interpret analytics feedback. While the dashboard visualizes the fragmented workflows and trial-and-error strategies described by Robinson et al.~\cite{robinson2022error}, recognizing the metacognitive value of these visualizations as advocated by Schwendimann et al.~\cite{schwendimann2016perceiving} likely requires longitudinal engagement beyond a 17-minute testing session. Qualitative feedback strongly supported this discrepancy between ease of use and perceived usefulness. Multiple participants expressed that the tool currently feels more aligned with an instructor's perspective, with one student noting, ``I feel like the dashboard is more of data analysis rather than for the student itself.'' Another participant explicitly requested actionable guidance, stating: ``I think it would be more useful for a teacher. I think it would be more useful for students to see how to solve a specific error, instead of just getting the feedback which errors you made.'' This directly echoes existing research on notebook-based analytics tools, which suggests that while fine-grained interaction data can reveal important workflow insights, learners may initially perceive such transparency as descriptive rather than transformative~\cite{cai2025jupyter,JELAI}. In other words, students may recognize that the dashboard shows them what they did, but may not yet connect this information to improved future performance without explicit instructional tips. Interestingly, despite these reservations about general usefulness, students did identify specific situational benefits. As one learner pointed out, ``I think this dashboard is useful when the learner wants to do a coding exam, so they know how much time they spent on one task,'' indicating that the perceived value of the tool may depend on the specific educational context.

\subsubsection{Overall User Experience and Workload}
Although not explicitly bound to a singular research question, assessing the subjective mental workload was a fundamental component in evaluating the dashboard's overall viability. A central and highly positive finding of this study is the low subjective mental workload measured via the NASA Task Load Index (NASA-TLX), which yielded a mean overall score of $M = 8.15$ ($SD = 2.92$) on a 20-point scale. When linearly transformed to a 100-point scale (approx. $40.75$), this value sits comfortably below the global median of $43.8$ reported in comprehensive user-experience meta-analyses~\cite{hertzum2021reference}.
This result directly addresses a critical concern in the design of educational tools. Computational notebooks already present a dense interactive structure by mixing narrative text, code, and visual output. Introducing an additional layer of fine-grained analytics inherently risks overwhelming the user. However, the low NASA-TLX score strongly suggests that the dashboard fulfills the critical requirement of being a low-workload addition to the programming environment. As emphasized by Corrin and De Barba~\cite{corrin2014exploring}, for learning analytics tools to be genuinely effective, they must be seamlessly integrated into the student's workflow without increasing cognitive burden.
The workload findings confirm that the dashboard successfully avoids this ``cognitive burden'' trap. It provides students with transparent, cell-level execution histories and visual timelines without distracting them from their primary goal of solving Python exercises. In combination with the high Perceived Ease of Use (PEOU) reported in RQ2, these results indicate that while students may still need pedagogical guidance to translate the dashboard's descriptive data into actionable steps, the system itself does not impose a significant cognitive hurdle. This lays a robust foundation for future iterations to focus on enhancing the instructional value of usefulness without needing to overhaul the underlying interaction paradigm.

\subsection{Limitations}
While this study provides valuable initial insights into the usability and acceptance of a learner-facing Jupyter dashboard, several limitations must be acknowledged. First, the sample size of 22 students restricts the generalizability of the findings and the statistical power necessary to detect significant correlations between perceived workload, usability, and technology acceptance. Although the sample size was sufficient for an exploratory evaluation and the scales demonstrated acceptable to excellent internal consistency, a larger and more diverse sample size would be required to draw robust, population-level conclusions. In addition to that, the short duration of the testing phase represents a significant constraint. Participants interacted with the dashboard while solving Python exercises for an average of approximately 17 minutes. As discussed regarding the moderate Perceived Usefulness scores, developing metacognitive skills and recognizing the instructional value of workflow reflection often requires long-term engagement and explicit pedagogical assistance~\cite{corrin2014exploring}. The brief exposure in this study captures initial impressions but may not accurately reflect how students would perceive, integrate, and benefit from the dashboard over the course of an entire academic semester.

\subsection{Future Work}
Building upon the foundational usability and technology acceptance findings of this exploratory study, future research should transition from subjective user evaluations to measuring objective pedagogical outcomes. A promising direction for subsequent research would be to conduct controlled experimental studies comparing learners with and without access to the dashboard over multiple learning units. Such studies could investigate whether providing visual feedback on cell execution histories, error rates, and time spent on tasks directly translates into improved programming performance. Researchers could focus on quantifying performance improvements through metrics such as the reduction in error rates between coding sessions, increased execution efficiency e.g., fewer retries per cell, and overall task correctness. To further refine the assessment of execution efficiency, future iterations of the system could utilize specific cell execution durations to evaluate whose code is more computationally optimized. Additionally, to accurately measure overall task correctness, the dashboard should integrate a database of verified solutions. Currently, code that executes without triggering an exception may be interpreted as successful, cross-referencing outputs with a validated database would allow the system to distinguish between code that merely executes without errors and code that is functionally correct.
Furthermore, to better understand the mechanisms of self-regulated learning, future iterations could integrate explicit reflection prompts directly into the dashboard interface. By combining system logs with pre- and post-surveys measuring variables like programming self-efficacy and motivation, researchers could determine whether increased metacognitive awareness actively mediates problem-solving success. 
Finally, while the current project focused exclusively on a student-facing dashboard, future work should expand this concept by developing and evaluating an instructor-facing counterpart. Gathering qualitative and quantitative feedback from educators will be essential to understand how aggregating this fine-grained student interaction data could facilitate targeted pedagogical interventions and reveal common coding difficulties across an entire cohort. Expanding these analytics to include direct comparisons between different students could provide both instructors and learners with valuable benchmarking insights, highlighting diverse problem-solving approaches and helping educators identify outliers who may need additional support or advanced challenges.